\documentclass[12pt, a4paper]{report}
\usepackage{mathrsfs, amsmath,amssymb,amsthm,mathtools,setspace,graphicx,array,tikz,url}
\usepackage{cite}
\usepackage[hidelinks]{hyperref}
\usepackage{fancyhdr}
\usepackage{lastpage} 
\usepackage[top=3cm, left=3cm, right=3cm, bottom=3cm]{geometry}
\usepackage{enumitem}
\usepackage{titlesec}
\usepackage{etoolbox}
\usepackage{amsthm} % Required for theorem environments
\usetikzlibrary{decorations.markings}
\patchcmd{\thebibliography}{\advance\leftmargin\labelsep}{\labelsep=.5em\leftmargin=2.5em\itemindent=0em\advance\leftmargin\labelsep}
\titlespacing*{\section}{0pt}{3em}{2em}
\titlespacing*{\chapter}{0pt}{0pt}{2em}
\titleformat{\chapter}[display]
  {\normalfont\huge\bfseries}
  {}
  {0pt}
  {}
\titleformat{\section}
  {\normalfont\large\bfseries} % format for the whole title
  {}                           % section number format
  {0pt}                        % spacing between number and title
  {}                           % code before the title text

\setstretch{1.1}
\pagestyle{fancy}
\fancyhf{} % Clear header and footer
\fancyfoot[C]{\footnotesize \thepage\ of \pageref{LastPage}}

\renewcommand{\headrulewidth}{0pt}
\renewcommand{\qedsymbol}{$\blacksquare$}

\newtheoremstyle{spaced}
  {15pt}   % Space above (e.g., 10pt)
  {20pt}   % Space below (e.g., 10pt)
  {\itshape} % Body font (e.g., italic)
  {}       % Indentation of the header
  {\bfseries} % Theorem header font (e.g., bold)
  {.}      % Punctuation after theorem header
  {.5em}   % Space after theorem header
  {}       % Theorem head spec (can be left empty)

\theoremstyle{spaced}
\newtheorem{conjecture}{Conjecture}[section]
\newtheorem{theorem}{Theorem}[section]
\newtheorem{definition}{Definition}[section]
\newtheorem{example}{Example}[section]
\newtheorem{lemma}{Lemma}[section]
\newtheorem{proposition}{Proposition}[section]
\newtheorem{corollary}{Corollary}[section]
\newtheorem{remark}{Remark}[section]
\newtheorem{namedlemma}{Negation Lemma}[section]
\newtheorem{exercise}{}
\newcommand{\setexercise}[1]{\setcounter{exercise}{#1-1}}
\pretocmd{\chapter}{\setcounter{exercise}{0}}{}{}
\renewcommand{\theHexercise}{\thechapter.\arabic{exercise}}

\title{Complex Analysis Exam Problems}
\author{
  Pic S. Infiniti (Second Edition)
}

\date{Summer 2026-2027}

\begin{document}
\hypersetup{pageanchor=false}
\maketitle

\tableofcontents

\hypersetup{pageanchor=true}
\chapter{Fall 2025}
\section{Part I}
\begin{exercise}
	Define $\gamma \colon [0,2\pi] \to \mathbb{C}$ by
	\[
		\gamma(t)=e^{it}.
	\]
	Compute
	\[
		\int_{\gamma} z^n\,dz
	\]
	for every integer $n$.
\end{exercise}

\begin{exercise}
	Let $\gamma$ be the closed polygon
	\[
		[1-i,\,1+i,\,-1+i,\,-1-i,\,1-i].
	\]
	Find
	\[
		\int_{\gamma}\frac{1}{z}\,dz.
	\]
\end{exercise}

\begin{exercise}
	\begin{enumerate}[(a)]
		\item Show that a M\"obius transformation has $0$ and $\infty$ as its only fixed points if and only if it is a dilation, but not the identity.

		\item Show that a M\"obius transformation has $\infty$ as its only fixed point if and only if it is a translation, but not the identity.
	\end{enumerate}
\end{exercise}

\begin{exercise}
	Let $G$ be a region, and suppose that $f\colon G\to\mathbb{C}$ is analytic such that $f(G)$ is a subset of a circle. Show that $f$ is constant.
\end{exercise}

\begin{exercise}
	Let
	\[
		f(z)=e^x\sin y+i e^x\cos y,
		\qquad z=x+iy.
	\]
	Determine whether $f$ is analytic on the complex plane.
\end{exercise}

\begin{exercise}
	Map
	\[
		G=\mathbb{C}\setminus\{z:-9\leq z\leq 9\}
	\]
	onto the open unit disk by an analytic function $f$. Can $f$ be one-to-one?
\end{exercise}

\begin{exercise}
	Give the principal branch of
	\[
		\sqrt{1-z^3}.
	\]
\end{exercise}
\section{Part II}
\begin{exercise}
	Compute the integral
	\[
		\int_{\gamma}
		\frac{e^z}{(z-2)^2(z-5)^5}\,dz,
	\]
	where $\gamma$ is the closed path
	\[
		\gamma(t)=3e^{it},
		\qquad 0\leq t\leq 2\pi.
	\]
\end{exercise}

\begin{exercise}
	Compute the integral in the previous exercise with $\gamma$ the closed path
	\[
		\gamma(t)=e^{it},
		\qquad 0\leq t\leq 2\pi.
	\]
\end{exercise}
\end{document}
